% Generated from validated Article Draft Result. Do not hand-edit; revise the structured draft instead.
\section{Detailed Chronology — 2024年後半の主要lifecycle}
\noindent\textbf{本文のテーマ別整理とは別に、model preview、release、product/API availability、研究発表を時系列で保持する。類似名称を一つの出来事へ畳み込まず、event identityを優先する。}\autocite{src-10b59b402b80,src-c3a08c57269d}

\begin{itemize}
  \item 2024年後半のchronologyでは、OpenAIのmodel/product/API系更新、FLUX系の公開、Amazon Nova、Googleの年末model/media更新、Llama familyの段階的更新、xAI Grok-2などを、それぞれ一次資料で確認できるeventとして保持する。preview、beta、GA、model release、API availability、research/publicationは別種のeventであり、同じfamily名でも後続更新へ吸収しない。
\end{itemize}
\autocite{src-10b59b402b80,src-1572648d0de9,src-5db07766b472,src-6bbe9721db06,src-99e6401ec31b,src-bc9a95ff1e00,src-c3a08c57269d,src-e92e95daf8e0}

この時系列を読む際は、同じvendorやmodel familyに属するという理由だけでeventを統合しない。特にpreviewから製品化、API提供、量子化版・minor family更新への移行は、利用者が実際に触れられるsurfaceやavailabilityが変わるため独立したchronology factとして残す。\autocite{src-10b59b402b80,src-6bbe9721db06}


\begin{claimboundary}
Chronologyは『何がいつ公開・提供されたか』を優先する。各発表に含まれるbenchmark、性能優位、品質表現はevent dateの事実とは別の帰属claimであり、時系列表に置かれたことを理由に独立検証済み事実へ格上げしない。\autocite{src-bc9a95ff1e00}
\end{claimboundary}
